\subsection{Collective Gravitational Coupling and Operational Geometry}
  \label{subsec:collective-coupling}

  Localized excitations act as persistent saturations of the local
  projective capacity, constraining the advancement of the admissible
  ordering in their neighborhood.
  Their collective influence is summarized by a response operator
  $K_{ij}$, the dressed counterpart of the bare connectivity
  $K_{0,\mathrm{bare}}$
  (Section~\ref{subsec:renormalization-hbar}).

  In the weak-structure regime, the effective potential satisfies
  \begin{equation}
    \nabla^2 \Phi_{\text{eff}}
    \approx 4\pi G_{\text{eff}}\,\rho,
    \label{eq:poisson-emergent}
  \end{equation}
  with $G_{\text{eff}}$ a phenomenological coupling.
  Earlier versions expressed it as
  $G_{\text{eff}} \approx c^4/(K_{0,\text{eff}}\,\chi_{c,\text{eff}}^2)$.
  That expression is withdrawn together with the relation between $K_0$, $\chi_c$ and $G$
  (Appendix~\ref{subsec:fundamental-quantities}); no expression of $G_{\text{eff}}$ in terms of substrate
  parameters is established in this paper.
  Spatial geometry is defined operationally: configurations are close if
  perturbations propagate efficiently between them.
  Gravity is recovered as a macroscopic manifestation of the local
  constraint on projective capacity imposed by matter configurations.
